\chapter{Interview-Grade Problems}
\label{ch:interview_problems}
\chaptermark{Interview-grade problems}

Quant interviews test probability, stochastic processes, linear
algebra, dynamic programming, and the ability to talk through all
four under pressure.  This chapter walks \textbf{fifteen problems
end-to-end} using a five-field template (problem statement,
first attempt, what goes wrong, correct approach, interviewer
variations) to expose the handful of patterns that recur across
Goldman Sachs and Jane Street questions.  The canonical industry
reference is \emph{A Practical Guide to Quantitative Finance
Interviews} by Xinfeng Zhou, the ``Green Book''
\citeGreenBook; we follow its conventions and cross-reference it
where a problem appears nearly verbatim.

\begin{backgroundbox}[Prerequisites]
The reader should be comfortable with:
\begin{itemize}
  \item \textbf{Probability} (\cref{ch:how_to_read}
    Appendix~A and \cref{ch:kyle_gm}):
    conditional expectation, Bayes' rule, independence, expectation
    and variance of Bernoulli, Binomial, Geometric, Poisson, Normal.
  \item \textbf{Stochastic processes} (\cref{ch:mean_reversion_ou}):
    the Brownian motion SDE $dW_t$, the Ornstein--Uhlenbeck process
    $dX_t = -\theta X_t\,dt + \sigma\,dW_t$, and martingales.
  \item \textbf{It\^o calculus and Black--Scholes}
    (\cref{ch:black_scholes}): It\^o's lemma on $f(W_t, t)$, the
    GBM SDE, and $C = S\CDFN(d_1) - K e^{-rT}\CDFN(d_2)$.
  \item \textbf{Linear algebra}: eigenvalues/eigenvectors, the
    spectral theorem for symmetric matrices, SVD, and that PCA
    diagonalises a covariance matrix.
  \item \textbf{Algorithms}: memoised recursion, the 2D DP table,
    and $O(nm)$ analysis.
\end{itemize}
Every result is derived here or pointed to its source chapter.
\end{backgroundbox}

Each problem uses the same five-field format; read them as
interview transcripts.

%% ================================================================
\section{Probability and combinatorics}
\label{sec:prob-problems}
%% ================================================================

%% ---------------- Problem 1: Monty Hall ----------------
\subsection{Problem 1: the Monty Hall door switch}
\label{sec:montyhall}

\paragraph{Problem statement.}
Three doors; behind one is a car, behind the other two are goats.
You pick door~1.  The host, who knows where the car is, opens
door~3 to reveal a goat.  He offers you the chance to switch to
door~2.  Should you switch?  What is the probability of winning
the car under each strategy?

\paragraph{First attempt.}
``Two doors remain, one has a car, so it is $1/2$ either way.''
This is the answer almost every candidate gives first.

\paragraph{What goes wrong.}
Treating the host's action as if it carried no information.  The
door the host opens is not random: he must open a goat door and
must not open yours.  That conditioning changes the posterior.

\paragraph{Correct approach.}
Condition on the initial pick.  With probability~$1/3$ you
initially picked the car; conditional on that, the host opens
either remaining door and switching loses.  With probability~$2/3$
you initially picked a goat; conditional on that, the host is
\emph{forced} to open the other goat door, leaving the car as the
door you could switch to.  Hence
\[
\Prob(\text{win} \mid \text{switch}) = \tfrac{2}{3},
\qquad
\Prob(\text{win} \mid \text{stay}) = \tfrac{1}{3}.
\]
A cleaner way: the initial pick is correct with probability $1/3$,
and the host's reveal is a deterministic function of your pick and
the true location, so the posterior on ``your pick is the car''
stays at $1/3$.  The remaining $2/3$ concentrates on the unopened
door you did not pick.

\paragraph{Interviewer variations.}
(i) Host picks a door uniformly and happens to reveal a goat: by
Bayes the posterior is $1/2$ on each remaining door.  (ii) $n$
doors, host opens $n-2$ goat doors: switching wins with
probability $(n-1)/n$.  (iii) Pick \emph{before} the host reveals
anything: both strategies give $1/2$.

\begin{intuitionbox}
The host is your accomplice.  Whenever you pick wrong (two thirds
of the time), he is forced to open a goat door and thereby
\emph{tells you which remaining door has the car}.  The paradox
dissolves once you see his action is constrained, not random.
\end{intuitionbox}

%% ---------------- Problem 2: Two children ----------------
\subsection{Problem 2: two children, at least one boy}
\label{sec:twokids}

\paragraph{Problem statement.}
A family has two children.  You are told that at least one of
them is a boy.  What is the probability that both are boys?
Separately: you are told that the \emph{elder} child is a boy.
What is the probability both are boys?

\paragraph{First attempt.}
``One child is a boy, so the other is independently a boy or girl,
giving $1/2$'' for both.

\paragraph{What goes wrong.}
The two questions condition on different events: treating them as
the same discards the difference between ``at least one'' and
``the elder is.''

\paragraph{Correct approach.}
The four equally likely birth orderings are
$\{BB, BG, GB, GG\}$.  Conditioning on ``at least one boy''
eliminates $GG$ and leaves three outcomes $\{BB, BG, GB\}$, of
which only $BB$ is two boys; the answer is $1/3$.  Conditioning
on ``the elder is a boy'' eliminates $GB$ and $GG$ and leaves
$\{BB, BG\}$; the answer is $1/2$.  The difference is that the
second piece of information distinguishes an ordered pair, while
the first does not.

\paragraph{Interviewer variations.}
(i) ``At least one boy born on Tuesday'': posterior $13/27
\approx 0.481$ (the ``Tuesday boy'' paradox); extra distinguishing
labels push toward $1/2$.  (ii) Three children, at least one boy:
$4/7$, by enumerating eight orderings minus $GGG$.

\begin{pitfallbox}[The ``unspecified child'' trap]
``One'' is ambiguous in English but precise in probability.  ``At
least one is a boy'' conditions on a set event; ``a specific one
(the elder) is a boy'' conditions on a coordinate.  On the
ambiguous version, pin down \emph{which} conditional is meant.
\end{pitfallbox}

%% ---------------- Problem 3: 100 prisoners ----------------
\subsection{Problem 3: the 100 prisoners and the boxes}
\label{sec:prisoners}

\paragraph{Problem statement.}
One hundred prisoners are numbered 1 to 100.  In a room are 100
closed boxes; inside each is a slip bearing one of the
numbers 1 to 100, with each number appearing exactly once in some
uniformly random permutation.  Each prisoner in turn enters the
room, opens at most 50 boxes, and must find the slip bearing
their own number.  They may not communicate during the game and
cannot leave evidence behind.  All 100 must succeed for the group
to win.  Find a strategy that wins with probability at least $30\%$.

\paragraph{First attempt.}
``Each prisoner opens 50 random boxes: individual success $1/2$,
group success $2^{-100}$.''  Correct under random opening, which
makes any $30\%$ bound sound impossible.

\paragraph{What goes wrong.}
The naive analysis assumes independence.  A coordinated strategy
correlates the searches so almost all succeed or almost all fail.

\paragraph{Correct approach.}
Prisoner $k$ opens box $k$; if the slip reads $j$, they open box
$j$, then the box labelled by its slip, forming a cycle through
the permutation.  Prisoner $k$ wins iff their cycle has length at
most 50; \emph{all} win iff the permutation has no cycle longer
than 50.

For a uniform random permutation of $n$ elements, the probability
of a cycle of length $\ell > n/2$ is exactly $1/\ell$: such a long
cycle is unique and placeable in $\binom{n}{\ell}(\ell-1)!(n-\ell)!$
ways out of $n!$.  Hence
\[
\Prob(\text{fail})
= \sum_{\ell=51}^{100} \frac{1}{\ell},
\qquad
\Prob(\text{win})
= 1 - \sum_{\ell=51}^{100} \frac{1}{\ell}
\approx 0.3118.
\]
As $n \to \infty$ the failure probability approaches
$\ln 2 \approx 0.693$, so the winning probability approaches
$1 - \ln 2 \approx 0.307$.

\paragraph{Interviewer variations.}
(i) Open 60 boxes each: failure $\sum_{\ell=61}^{100} 1/\ell
\approx 0.502$, win $\approx 0.498$.  (ii) Can you beat
cycle-following?  No: an information-theoretic argument.  (iii)
Verify $1/\ell$ by Monte Carlo.

\begin{keyfactbox}[Cycle lemma for random permutations]
For a uniform random permutation of $n$ symbols and any
$\ell > n/2$, the probability the permutation contains a cycle of
length exactly $\ell$ is $1/\ell$.  The cycle-following
strategy in the 100-prisoner problem therefore wins with probability
$1 - \sum_{\ell=51}^{100} 1/\ell \approx 0.3118$.
\end{keyfactbox}

%% ---------------- Problem 4: Secretary / optimal stopping ----------------
\subsection{Problem 4: the secretary problem}
\label{sec:secretary}

\paragraph{Problem statement.}
$n$ candidates are interviewed in uniformly random order.  After
each interview you must either hire the candidate on the spot or
dismiss them forever.  You can rank any two you have seen but not
candidates you have not.  What strategy maximises the probability
of hiring the single best candidate?

\paragraph{First attempt.}
``Hire the first candidate better than everyone so far.''  Right
in spirit, but leaves open when to \emph{start} applying the rule.

\paragraph{What goes wrong.}
The rule triggers too early: the best-so-far among the first few
is not a strong signal about the global best.  A \emph{warm-up}
window is needed.

\paragraph{Correct approach.}
Fix threshold $r \in \{1, \dots, n\}$.  Interview the first $r-1$
(hire none), then hire the first candidate better than everyone so
far.  The success probability is
\[
P(r)
= \frac{r-1}{n} \sum_{i=r}^{n} \frac{1}{i-1}.
\]
$(r-1)/n$ is the probability the best appears after position
$r-1$; the sum accounts for the best-so-far at position $i-1$
lying in the warm-up.  Write $P(r) \approx (r/n) \ln(n/r)$ for
large $n$ (harmonic sum as an integral); setting $dP/dr = 0$
gives $\ln(n/r) = 1$, i.e.\ $r = n/e$, and $P \to 1/e \approx
0.3679$.

For $n = 10$ the optimum is $r = 3$ (reject the first two, then
hire the next new best), success $\approx 0.3987$.

\paragraph{Interviewer variations.}
(i) Minimise \emph{expected rank} instead: limiting value
$\approx 3.87$.  (ii) ``Postdoc'' variant: hire the \emph{second}
best.  (iii) Two hires, a staple follow-up.

\begin{historybox}[title={The 1/e law}]
The secretary problem and its $1/e$ threshold date from the 1950s;
Martin Gardner popularised it in his February 1960 \emph{Scientific
American} column as ``the Googol game.''  It reappears in
interviews because it teaches \emph{optimal stopping under no
prior}, which generalises to bidding, exercise, and liquidation.
\end{historybox}

%% ---------------- Problem 5: Coupon collector ----------------
\subsection{Problem 5: the coupon collector}
\label{sec:coupon}

\paragraph{Problem statement.}
Cereal boxes contain coupons, each of which is independently one
of $n$ equally likely types.  How many boxes do you expect to
open before you have collected at least one of every type?

\paragraph{First attempt.}
``With $n$ types I expect roughly $n$ boxes.''

\paragraph{What goes wrong.}
The first coupon is certain to be new, but the last is rare.  As
the collection nears completion the waiting time stretches out:
the $k$-th new coupon, given $k-1$ already collected, is geometric
with success probability $(n-k+1)/n$.  Expected waiting time is
$n/(n-k+1)$.

\paragraph{Correct approach.}
Summing over $k = 1, \dots, n$,
\[
\E[T_n]
= \sum_{k=1}^{n} \frac{n}{n-k+1}
= n \sum_{j=1}^{n} \frac{1}{j}
= n H_n,
\]
where $H_n$ is the $n$-th harmonic number.  For large $n$ the
expected time is $n \ln n + \gamma n + O(1)$ with
$\gamma \approx 0.5772$ the Euler--Mascheroni constant.  For a
deck of $n = 52$ cards, $\E[T] \approx 235.98$.

The variance is
\(
\Var(T_n) = n^2 \sum_{k=1}^{n} 1/k^2 - n H_n
\sim \pi^2 n^2 / 6,
\)
so the standard deviation scales linearly in $n$, of the same
order as the mean.

\paragraph{Interviewer variations.}
(i) Unequal coupon probabilities: expected time grows; use
inclusion--exclusion.  (ii) Variance derivation: decompose $T_n$
as a sum of independent geometrics and sum variances.  (iii)
Concentration: Chernoff-style bounds give $\Prob(T_n > n \ln n +
cn) \sim e^{-c}$.

\begin{intuitionbox}
The harmonic sum is where the time goes.  The first half of types
costs $\approx n \ln 2 \approx 0.69 n$; the last few cost another
$\approx 0.5 n \ln n$.  The tail dominates.
\end{intuitionbox}

%% ================================================================
\section{Stochastic processes}
\label{sec:stoch-problems}
%% ================================================================

%% ---------------- Problem 6: Brownian first hitting time ----------------
\subsection{Problem 6: first hitting time of a Brownian motion}
\label{sec:hitting}

\paragraph{Problem statement.}
Let $W_t$ be a standard Brownian motion starting at $W_0 = 0$.
Let $T_a = \inf\{t \ge 0 : W_t = a\}$ for $a > 0$.  What is
$\Prob(T_a < \infty)$?  What is $\E[T_a]$?

\paragraph{First attempt.}
``The Brownian motion is a martingale with mean zero; if I apply
optional stopping at $T_a$ I get $\E[W_{T_a}] = 0$, but
$W_{T_a} = a$, contradiction.  So $T_a$ must be infinite with
positive probability.''

\paragraph{What goes wrong.}
The optional-stopping theorem has hypotheses.  Applying it to an
unbounded stopping time without checking uniform integrability is
the classic mistake here.

\paragraph{Correct approach.}
\emph{Step 1: hitting is certain.}  By recurrence of
one-dimensional Brownian motion, $\Prob(T_a < \infty) = 1$ for
any finite $a$.  Equivalently, $W_t$ is recurrent in dimension 1
so every level is visited.  \emph{Step 2: mean is infinite.}
Apply It\^o's lemma to $e^{\lambda W_t - \lambda^2 t/2}$ to obtain
a martingale, evaluate at $T_a$, and use optional stopping on the
bounded stopping time $T_a \wedge N$.  Taking $N \to \infty$
yields
\[
\E[e^{-\lambda^2 T_a/2}] = e^{-\lambda a}, \qquad \lambda > 0,
\]
the Laplace transform of the hitting-time density.  Inverting,
the density is
\[
f_{T_a}(t) = \frac{a}{\sqrt{2\pi t^3}}\, e^{-a^2/(2t)},
\qquad t > 0,
\]
a \emph{L\'evy distribution}.  Its mean is
$\int_0^\infty t f_{T_a}(t)\,dt = \infty$ because of the heavy
$t^{-3/2}$ tail.  So $T_a$ is almost surely finite but has
infinite expectation.  \emph{Step 3: scaling.}  By the scaling
$W_t \stackrel{d}{=} a^{-1} W_{a^2 t}$, one reads off
$T_a \stackrel{d}{=} a^2 T_1$: first passage scales with the
square of the level.  Dimensional analysis alone gives this:
$[a^2/t]$ has units of volatility squared.

\paragraph{Interviewer variations.}
(i) Hit either $a$ or $-b$ first with $a, b > 0$: the probability
of hitting $a$ first is $b/(a+b)$ and the expected time is $ab$.
(ii) Add a drift $\mu$: the Laplace transform picks up a drift
term and the hitting distribution becomes inverse Gaussian.
(iii) Prove the reflection principle and use it to derive
$\Prob(\max_{s \le t} W_s \ge a) = 2\,\Prob(W_t \ge a)$.

\begin{keyfactbox}[Brownian hitting time]
For $W_0 = 0$, $a > 0$:
\[
\Prob(T_a < \infty) = 1, \quad
\E[T_a] = \infty, \quad
T_a \stackrel{d}{=} a^2 T_1,
\]
with density $f_{T_a}(t) = \frac{a}{\sqrt{2\pi t^3}} e^{-a^2/(2t)}$.
Finite in probability, infinite in mean.
\end{keyfactbox}

%% ---------------- Problem 7: OU half-life ----------------
\subsection{Problem 7: half-life of an OU process}
\label{sec:ou-half}

\paragraph{Problem statement.}
Let $X_t$ satisfy the Ornstein--Uhlenbeck SDE
$dX_t = -\theta (X_t - \mu)\,dt + \sigma\,dW_t$ with
$\theta > 0$.  Derive the \emph{half-life} of mean reversion:
the expected time for a displacement from $\mu$ to decay to half
its initial value.

\paragraph{First attempt.}
``Half-lives are $1/\theta$.''

\paragraph{What goes wrong.}
The time constant is $1/\theta$; the half-life differs by a
factor of $\ln 2$.  Confusing the two costs full credit.

\paragraph{Correct approach.}
Take expectations of the SDE; the noise term is zero-mean, so
\[
\frac{d\,\E[X_t]}{dt}
= -\theta \left(\E[X_t] - \mu\right),
\]
with solution $\E[X_t] - \mu = (X_0 - \mu) e^{-\theta t}$.  The
half-life $t_{1/2}$ satisfies $e^{-\theta t_{1/2}} = 1/2$:
\[
t_{1/2} = \frac{\ln 2}{\theta} \approx \frac{0.693}{\theta}.
\]
For $\theta = 0.1$/day, $t_{1/2} \approx 6.93$ days.

\paragraph{Interviewer variations.}
(i) Stationary distribution: Gaussian, mean $\mu$, variance
$\sigma^2/(2\theta)$.  (ii) Pair-trade holding horizon with
residual half-life $t_{1/2}$: a few half-lives.  (iii) Discrete
time: $\rho = e^{-\theta \Delta t}$, so $t_{1/2} = -\ln 2 \cdot
\Delta t / \ln \rho$.

\begin{intuitionbox}
``Time constant'' is when the deviation falls to $1/e \approx
0.368$; ``half-life'' is when it falls to $0.5$.  The two differ
by $\ln 2 \approx 0.693$.  Use half-lives with traders; use time
constants in the calculus.
\end{intuitionbox}

%% ---------------- Problem 8: Gambler's ruin ----------------
\subsection{Problem 8: gambler's ruin and martingale optional stopping}
\label{sec:ruin}

\paragraph{Problem statement.}
Start with $a$ chips at a fair game; each round win or lose one
chip with probability $1/2$.  The game ends when the stack hits
either $0$ or $N$.  What is the probability of exit at $N$?
What is the expected number of rounds until exit?

\paragraph{First attempt.}
``Random walk, so maybe $a/N$?''  The guess is correct; justifying
it earns the credit.

\paragraph{What goes wrong.}
The ``distance-weighted'' argument leaves two holes: why is it
exactly linear in $a$, and why may we use optional stopping?

\paragraph{Correct approach.}
\emph{Probability of reaching $N$.}  Let $X_t$ denote the chip
stack.  $X_t$ is a martingale because the game is fair, so by
optional stopping (the stopping time $\tau$ is bounded in
expectation and the martingale is bounded)
$\E[X_\tau] = X_0 = a$.  Writing
$\E[X_\tau] = 0 \cdot \Prob(\text{ruin}) + N \cdot \Prob(\text{hit } N)$
gives
\[
\Prob(\text{hit } N) = \frac{a}{N}.
\]
\emph{Expected duration.}  Apply optional stopping to the second
martingale $X_t^2 - t$.  This is a martingale because
$\E[X_{t+1}^2 - X_t^2 \mid X_t] = 1$ for a fair $\pm 1$ random
walk.  So $\E[X_\tau^2] = \E[X_0^2] + \E[\tau]$.
\[
\E[X_\tau^2]
= 0 \cdot \Prob(\text{ruin}) + N^2 \cdot \frac{a}{N}
= a N,
\]
so $\E[\tau] = aN - a^2 = a(N-a)$.  With $a = 10$, $N = 100$:
$\Prob(\text{win}) = 10\%$, $\E[\tau] = 900$ rounds.

\emph{Unfair coin.}  If the win probability is $p \ne 1/2$ with
$q = 1-p$ and $r = q/p$, the martingale $r^{X_t}$ gives
$\Prob(\text{hit } N) = (1-r^a)/(1-r^N)$.  Taking $r \to 1$
recovers $a/N$.

\paragraph{Interviewer variations.}
(i) Asymmetric barriers $[-b, a]$ in a Brownian setting.  (ii)
Why finite expected duration?  Recurrent walk, absorbing
barriers.  (iii) Kelly connection: in a biased game, the
gambler's-ruin formula justifies sizing below Kelly to control
ruin probability (\cref{ch:kelly_risk}).

%% ---------------- Problem 9: Itô on $W^2$ and $e^W$ ----------------
\subsection{Problem 9: It\^o's lemma on $W_t^2$ and $e^{W_t}$}
\label{sec:ito}

\paragraph{Problem statement.}
Use It\^o's lemma to derive the SDEs for $W_t^2$ and $e^{W_t}$,
and identify an explicit martingale built from each.

\paragraph{First attempt.}
``By the chain rule, $d(W^2) = 2W\,dW$ and $d(e^W) = e^W\,dW$.''

\paragraph{What goes wrong.}
That is the ordinary chain rule.  It\^o's lemma adds
$\tfrac{1}{2} f''(W_t)\,dt$ from the quadratic variation of
Brownian motion, since $(dW)^2 = dt$.  Forgetting this is the
definitional error of stochastic calculus.

\paragraph{Correct approach.}
Apply It\^o's lemma $df(W_t, t) = f_t\,dt + f_W\,dW_t
+ \tfrac{1}{2} f_{WW}\,dt$.

\emph{Case 1: $f(W) = W^2$.}  Here $f_W = 2W$, $f_{WW} = 2$, so
\[
d(W_t^2) = 2 W_t\,dW_t + dt.
\]
Integrating, $W_t^2 - t = 2\int_0^t W_s\,dW_s$, the
stochastic-integral form of the L\'evy martingale.  Check:
$\E[W_t^2 - t] = 0$, and increments $W_t^2 - W_s^2 - (t-s)$ have
zero conditional mean given $\mathcal F_s$.

\emph{Case 2: $f(W) = e^W$.}  Here $f_W = f_{WW} = e^W$, so
\[
d(e^{W_t}) = e^{W_t}\,dW_t + \tfrac{1}{2} e^{W_t}\,dt.
\]
The drift is nonzero, so $e^{W_t}$ is \emph{not} a martingale.
To kill the drift, multiply by $e^{-t/2}$: It\^o's product rule
gives
\[
d(e^{W_t - t/2}) = e^{W_t - t/2}\,dW_t,
\]
which is a pure stochastic integral and therefore a martingale.
This is the \emph{exponential martingale}, the building block of
risk-neutral pricing.

\paragraph{Interviewer variations.}
(i) Apply It\^o to $W_t^3$: $d(W^3) = 3W^2 dW + 3W\,dt$, and
$W_t^3 - 3t W_t$ is a martingale.  (ii) For geometric Brownian
motion $S_t = S_0 e^{(\mu - \sigma^2/2)t + \sigma W_t}$, show
$dS_t = \mu S_t\,dt + \sigma S_t\,dW_t$.  (iii) Use the
exponential martingale as the Radon--Nikodym density in
Girsanov's theorem.

\begin{keyfactbox}[It\^o lemma and two canonical martingales]
\[
W_t^2 - t \quad\text{and}\quad e^{W_t - t/2}
\]
are martingales: the first from $d(W^2) = 2W\,dW + dt$, the
second from the drift correction $\tfrac{1}{2} e^W\,dt$.  These
two drive every derivation in \cref{ch:black_scholes}.
\end{keyfactbox}

%% ================================================================
\section{Linear algebra}
\label{sec:linalg-problems}
%% ================================================================

%% ---------------- Problem 10: PCA interview ----------------
\subsection{Problem 10: the PCA question}
\label{sec:pca}

\paragraph{Problem statement.}
You have $n$ vectors of returns $x_1, \dots, x_n \in \R^d$, one
per day, stacked into a matrix $X \in \R^{n \times d}$.  Explain
what principal component analysis does, derive the first
principal component, and say what the eigenvalues mean.

\paragraph{First attempt.}
``PCA finds the top eigenvectors of the covariance matrix.''
True but not a derivation.

\paragraph{What goes wrong.}
Quoting the answer fails the test.  The real question: how do
you go from \emph{maximise variance} to \emph{top eigenvector}?

\paragraph{Correct approach.}
Assume data are centred: column means of $X$ are zero.  The
sample covariance matrix is
$\Sigma = \tfrac{1}{n-1} X^\top X$.  We seek a unit vector
$w \in \R^d$ that maximises the variance of the projection
$X w$:
\[
\max_{w} \; w^\top \Sigma w \quad\text{subject to}\quad \|w\| = 1.
\]
Lagrangian: $\mathcal L = w^\top \Sigma w - \lambda(w^\top w - 1)$.
Setting $\nabla_w \mathcal L = 0$ gives $\Sigma w = \lambda w$,
so $w$ is an eigenvector of $\Sigma$ and $\lambda$ is the
eigenvalue.  Among all eigenvectors, the one maximising
$w^\top \Sigma w = \lambda$ is the top one.  Subsequent components
are obtained by maximising variance subject to orthogonality to
all prior components, which yields the next eigenvector in order.

The $k$-th eigenvalue is the variance along the $k$-th principal
component; total variance is $\mathrm{tr}(\Sigma) = \sum_k
\lambda_k$, so the variance explained by the first $K$ components
is $\sum_{k=1}^K \lambda_k / \sum_{k=1}^d \lambda_k$.
Numerically, PCA is done via SVD $X = U S V^\top$: principal
directions are columns of $V$, squared singular values are scaled
eigenvalues.

\paragraph{Interviewer variations.}
(i) Different feature scales?  Standardise first, or large-unit
features dominate.  (ii) When does PCA fail?  When variance is a
poor proxy for signal: rare important risk factors may have
low variance, noise may have high variance without predictive
content.  (iii) Factor models: top PCs of a rates curve are
level, slope, curvature (\cref{ch:rates_derivatives}).

%% ---------------- Problem 11: rank-1 update ----------------
\subsection{Problem 11: eigenvalues of a rank-1 update}
\label{sec:rank1}

\paragraph{Problem statement.}
Let $u \in \R^n$ be a nonzero vector.  What are the eigenvalues of
the $n \times n$ matrix $I + u u^\top$?

\paragraph{First attempt.}
``Expand the characteristic polynomial.''  Feasible for small $n$
but does not generalise.

\paragraph{What goes wrong.}
Determinant expansion is $O(n!)$ and hides the structure.  The
problem is trivial once you notice the 2D structure.

\paragraph{Correct approach.}
Split $\R^n$ into $\mathrm{span}(u)$ and its orthogonal
complement $u^\perp$.  For any $v \in u^\perp$,
\(
(I + u u^\top) v = v,
\)
so every vector in the $(n-1)$-dimensional space $u^\perp$ is an
eigenvector with eigenvalue $1$.  For $v = u$,
\(
(I + u u^\top) u = u + u (u^\top u) = (1 + \|u\|^2) u,
\)
so $u$ is an eigenvector with eigenvalue $1 + \|u\|^2$.  The
spectrum is therefore
\[
\{\underbrace{1, 1, \dots, 1}_{n-1}, 1 + \|u\|^2\}.
\]
Check the trace: $\sum \lambda_i = (n-1) + (1 + \|u\|^2)
= n + \|u\|^2 = \mathrm{tr}(I) + u^\top u = \mathrm{tr}(I + uu^\top)$. \checkmark

For $u = (1,2,3)^\top$, $\|u\|^2 = 14$, eigenvalues are
$\{1, 1, 15\}$, sum $17$ equals trace.

\paragraph{Interviewer variations.}
(i) ``Eigenvalues of $A + uv^\top$ for general $u, v$ (rank-1
update of general $A$)?''  Use the matrix determinant lemma
$\det(A + uv^\top) = \det(A)(1 + v^\top A^{-1} u)$.  (ii)
``Sherman--Morrison formula for $(A + uv^\top)^{-1}$?''
\(
(A + uv^\top)^{-1}
= A^{-1} - \tfrac{A^{-1} u v^\top A^{-1}}{1 + v^\top A^{-1} u}.
\)
(iii) ``How does this relate to Kalman filters?''  Kalman updates
covariance by rank-1 additions; Sherman--Morrison is the
$O(n^2)$ update path.

\begin{intuitionbox}
$I + uu^\top$ stretches by $1 + \|u\|^2$ along $u$ and leaves
$u^\perp$ unchanged.  See the geometry and the algebra writes
itself; this is the template for every low-rank update problem.
\end{intuitionbox}

%% ================================================================
\section{Dynamic programming}
\label{sec:dp-problems}
%% ================================================================

%% ---------------- Problem 12: LCS ----------------
\subsection{Problem 12: longest common subsequence}
\label{sec:lcs}

\paragraph{Problem statement.}
Given two strings $s$ of length $m$ and $t$ of length $n$, compute
the length of their longest common subsequence (not substring).
A subsequence preserves order but allows gaps.

\paragraph{First attempt.}
``Try every subset of $s$ and test for occurrence in $t$.''
$O(2^m n)$; infeasible for $m = 50$.

\paragraph{What goes wrong.}
Exponential blow-up ignores the subproblem structure: the LCS of
two prefixes depends only on LCSes of slightly shorter prefixes.

\paragraph{Correct approach.}
Let $L(i, j)$ denote the LCS length of $s[1..i]$ and $t[1..j]$.
The recurrence is
\[
L(i, j) =
\begin{cases}
0 & \text{if } i = 0 \text{ or } j = 0, \\
L(i-1, j-1) + 1 & \text{if } s_i = t_j, \\
\max(L(i-1, j), L(i, j-1)) & \text{otherwise}.
\end{cases}
\]
Filling the table bottom-up requires $O(nm)$ time and space.
Reconstructing the actual subsequence requires a backtrace from
$(m,n)$ to $(0,0)$.

For \texttt{ABCBDAB} and \texttt{BDCAB}, LCS length is 4 (e.g.,
\texttt{BCAB}).  The Python below computes it in eleven lines.

\begin{lstlisting}[caption={Longest common subsequence by dynamic programming.}]
def lcs(s, t):
    m, n = len(s), len(t)
    dp = [[0] * (n + 1) for _ in range(m + 1)]
    for i in range(1, m + 1):
        for j in range(1, n + 1):
            if s[i-1] == t[j-1]:
                dp[i][j] = dp[i-1][j-1] + 1
            else:
                dp[i][j] = max(dp[i-1][j], dp[i][j-1])
    return dp[m][n]

print(lcs("ABCBDAB", "BDCAB"))   # 4
\end{lstlisting}

\paragraph{Interviewer variations.}
(i) Reduce space to $O(\min(m,n))$ by keeping two rows.  (ii)
Longest \emph{common substring}: reset to zero on mismatch, answer
is max cell, not $L(m,n)$.  (iii) Edit distance (Levenshtein): 3
operations, similar $O(nm)$ DP.  (iv) LCS on compressed inputs
with $O(\log)$-alphabet is $NP$-hard.

%% ---------------- Problem 13: knapsack ----------------
\subsection{Problem 13: 0/1 knapsack}
\label{sec:knapsack}

\paragraph{Problem statement.}
Given items with weights $w_i$ and values $v_i$ for
$i = 1, \dots, n$, and a knapsack of capacity $W$, choose a
subset maximising total value subject to total weight
$\le W$.  Each item is either taken or left (no fractions).

\paragraph{First attempt.}
``Take items in decreasing value per unit weight.''  This is the
\emph{fractional} knapsack greedy; for 0/1 it is a heuristic,
not optimal.

\paragraph{What goes wrong.}
Greedy by ratio may miss a combination that trades off smaller,
worse-ratio items for a better fit.  Counter-example: weights
$[2, 3, 4, 5]$, values $[3, 4, 5, 6]$, capacity $W = 5$.  Greedy
by ratio picks item~1 (ratio 1.5) then item~3 (ratio 1.25) for
weight 6, which exceeds capacity.  The optimum picks items 1 and
2 (weight 5, value 7).

\paragraph{Correct approach.}
Let $K(i, w)$ denote the maximum value using a subset of the
first $i$ items with total weight $\le w$.  Recurrence:
\[
K(i, w) =
\begin{cases}
K(i-1, w) & \text{if } w_i > w, \\
\max\bigl(K(i-1, w),\; K(i-1, w - w_i) + v_i\bigr) & \text{otherwise}.
\end{cases}
\]
Base case $K(0, w) = 0$.  Answer $K(n, W)$.  Complexity is
$O(nW)$ time and space, which is \emph{pseudopolynomial} because $W$ is
encoded in $\log W$ bits; the problem is $NP$-hard when $W$ is
allowed to grow exponentially in the input size.

\begin{lstlisting}[caption={0/1 knapsack by dynamic programming.}]
def knapsack(weights, values, W):
    n = len(weights)
    dp = [[0] * (W + 1) for _ in range(n + 1)]
    for i in range(1, n + 1):
        for w in range(W + 1):
            dp[i][w] = dp[i-1][w]
            if weights[i-1] <= w:
                take = dp[i-1][w - weights[i-1]] + values[i-1]
                dp[i][w] = max(dp[i][w], take)
    return dp[n][W]

print(knapsack([2,3,4,5], [3,4,5,6], 5))   # 7
\end{lstlisting}

\paragraph{Interviewer variations.}
(i) Unbounded knapsack: $K(i,w) = \max(K(i-1,w), K(i,w-w_i)+v_i)$.
(ii) Fractional knapsack: greedy is optimal.  (iii) Subset-sum:
a special case.  (iv) FPTAS via scaled values.

\begin{pitfallbox}[title={Pseudopolynomial is not polynomial}]
$O(nW)$ sounds polynomial but is exponential in the input
bit-length since $W$ is written in $\log W$ bits.  For $n = 100$,
$W = 10^{12}$, the table has $10^{14}$ cells.  Under large $W$:
branch-and-bound, FPTAS, or a problem-specific reformulation.
\end{pitfallbox}

%% ================================================================
\section{Options and market intuition}
\label{sec:options-problems}
%% ================================================================

%% ---------------- Problem 14: Delta = N(d1) ----------------
\subsection{Problem 14: derive $\Delta = \CDFN(d_1)$}
\label{sec:deltad1}

\paragraph{Problem statement.}
Under Black--Scholes, the price of a European call is
$C = S \CDFN(d_1) - K e^{-rT} \CDFN(d_2)$ where
$d_{1,2} = [\ln(S/K) + (r \pm \sigma^2/2) T] / (\sigma \sqrt T)$.
Derive the delta $\Delta = \partial C / \partial S$.

\paragraph{First attempt.}
``$C$ is linear in $S$ through the first term, so
$\Delta = \CDFN(d_1)$.''  Right answer, incomplete derivation:
the $\CDFN$ arguments also depend on $S$.

\paragraph{What goes wrong.}
Candidates miss that $d_1, d_2$ are functions of $S$.  The result
survives only via a non-trivial cancellation.

\paragraph{Correct approach.}
Both $d_1$ and $d_2$ depend on $S$ through $\ln S$ with
$\partial d_1/\partial S = \partial d_2/\partial S = 1/(S \sigma
\sqrt T)$.  Differentiating,
\begin{align*}
\frac{\partial C}{\partial S}
&= \CDFN(d_1)
+ S \pdfN(d_1) \cdot \frac{1}{S \sigma \sqrt T}
- K e^{-rT} \pdfN(d_2) \cdot \frac{1}{S \sigma \sqrt T} \\
&= \CDFN(d_1) + \frac{1}{\sigma \sqrt T}
\left[\pdfN(d_1) - \frac{K e^{-rT}}{S} \pdfN(d_2) \right].
\end{align*}
The square bracket vanishes by the identity
\(
S \pdfN(d_1) = K e^{-rT} \pdfN(d_2),
\)
which is a direct consequence of
$d_1^2 - d_2^2 = 2 \ln(S/K) + 2 r T$ and the form of the
standard normal density.  (Plugging $\pdfN(x) = e^{-x^2/2}/\sqrt{2\pi}$
into the ratio $\pdfN(d_1)/\pdfN(d_2) = e^{-(d_1^2-d_2^2)/2}
= e^{-\ln(S/K) - rT} = (K/S) e^{-rT}$ gives the identity directly.)
So $\Delta = \CDFN(d_1)$.

Numerical check: $S = K = 100$, $r = 5\%$, $\sigma = 20\%$,
$T = 1$.  $d_1 = (0.05 + 0.02)/0.2 = 0.35$,
$\CDFN(0.35) \approx 0.6368$.  Call price $\approx 10.45$.

\paragraph{Interviewer variations.}
(i) ``Derive Gamma.''  Answer $\Gamma = \pdfN(d_1)/(S \sigma
\sqrt T)$.  (ii) ``Vega.''  $\mathrm{Vega} = S \pdfN(d_1)
\sqrt T$.  (iii) ``Put delta.''  $\Delta_P = \CDFN(d_1) - 1$
via put--call parity.  (iv) ``Why is delta also the risk-neutral
probability of expiring in the money?''  It is not exactly:
that probability is $\CDFN(d_2)$.  $\CDFN(d_1)$ is the probability in a
different measure (the stock numeraire).  Confusing $d_1$ and
$d_2$ here is a common pitfall.

\begin{pitfallbox}[title={Delta is not exercise probability}]
$\CDFN(d_2)$ is the risk-neutral probability the call finishes
in the money.  $\CDFN(d_1) = \Delta$ is the same probability under
the stock-as-numeraire measure.  They differ by $\sigma \sqrt T$
in the argument.  Interviewers use this to test numeraire change,
not formula recall.
\end{pitfallbox}

%% ---------------- Problem 15: Straddle P&L ----------------
\subsection{Problem 15: straddle P\&L under an implied-realised mismatch}
\label{sec:straddle}

\paragraph{Problem statement.}
An ATM straddle (long one call + one put, both struck at the spot)
on a $S = 100$ underlying has one year to expiry.  You pay
Black--Scholes prices using an implied volatility of $\sigma_{\mathrm{imp}}
= 20\%$.  The realised volatility over the next year turns out to
be $\sigma_{\mathrm{real}} = 80\%$.  Estimate the expected P\&L.

\paragraph{First attempt.}
``Realised vol is four times implied, so payoff is four times
cost.''  The ratio is roughly right but cost and payoff sit under
different functions of $\sigma$.

\paragraph{What goes wrong.}
Under GBM, $\E|S_T - S_0| \approx S_0 \sigma \sqrt{T} \sqrt{2/\pi}$,
linear in realised volatility, while the straddle price is
$\approx S_0 \sigma_{\mathrm{imp}} \sqrt{T} \sqrt{2/\pi}$, also
linear.  The payoff scales as $\sigma_{\mathrm{real}}/
\sigma_{\mathrm{imp}}$ but P\&L is payoff minus cost.

\paragraph{Correct approach.}
Under GBM with zero rates and zero drift for simplicity, the
expected straddle payoff is
\[
\E[|S_T - K|]
= S_0 \sigma_{\mathrm{real}} \sqrt{T/(2\pi)} \cdot 2
= S_0 \sigma_{\mathrm{real}} \sqrt{T} \sqrt{2/\pi}
\]
(the factor 2 comes from either side of the strike).  The ATM
straddle price at implied vol $\sigma_{\mathrm{imp}}$ is
approximately $2 \cdot S_0 \sigma_{\mathrm{imp}} \sqrt{T/(2\pi)}
= S_0 \sigma_{\mathrm{imp}} \sqrt{T} \sqrt{2/\pi}$ (using the
ATM approximation $C \approx 0.4 S_0 \sigma \sqrt{T}$).
Expected P\&L is therefore
\[
\E[\mathrm{P\&L}] \approx S_0 \sqrt{T} \sqrt{2/\pi}
(\sigma_{\mathrm{real}} - \sigma_{\mathrm{imp}}).
\]
With $S_0 = 100$, $T = 1$, $\sigma_{\mathrm{real}} = 0.8$,
$\sigma_{\mathrm{imp}} = 0.2$:
$100 \cdot 0.7979 \cdot 0.6 \approx 47.87$.  Expected payoff
$\approx 63.8$, cost $\approx 15.96$, P\&L $\approx 47.9$ per
straddle.  A 4$\times$ vol mismatch gives a roughly 4$\times$
P\&L because the cost is a small fraction of the realised payoff.

\paragraph{Interviewer variations.}
(i) ``What if realised is only 30\%?''  P\&L $\approx 7.98$, still
positive but modest.  (ii) ``What is the \emph{variance} of P\&L?''
Large: the straddle payoff is heavy-tailed; in 2020 March VIX
spikes the ratio exceeded 10x for weekly straddles.  (iii) ``Why
would anyone sell a straddle?''  The seller is short gamma and
collects theta; profits when realised < implied.  (iv) ``Delta
hedge the straddle: what is the P\&L?''  The classic ``gamma
P\&L'' identity: daily P\&L $\approx \tfrac{1}{2} \Gamma S^2
[(\Delta S/S)^2 - \sigma_{\mathrm{imp}}^2 \Delta t]$, summed over
the life of the option.

\begin{intuitionbox}
Option P\&L is the integral of squared returns minus squared
implied.  Realised 80\% vs implied 20\% is a 16x difference in
squared volatility, enormous for a long straddle.  The
intuition ``vol is mispriced by 4x'' is right; gamma P\&L accrues
at the difference of \emph{squared} vols.
\end{intuitionbox}

%% ================================================================
\section{Meta-patterns across the fifteen problems}
\label{sec:meta}
%% ================================================================

The problems span probability, stochastic processes, linear
algebra, DP, and pricing but share four moves.  Re-read this
section after a first pass: the pattern names are the vocabulary
you will use to think aloud.

\paragraph{Pattern 1: change of variables and dimensional analysis.}
Pass to log variables for GBM (Problem 9), to squared levels for
variance martingales (Problem 8), to units of $a^2$ or $\sigma^2
T$ for hitting times (Problem 6).  The right change of variables
linearises a nonlinear problem.

\paragraph{Pattern 2: look for a martingale first.}
Gambler's ruin (Problem~8) and Brownian hitting times (Problem~6)
reduce to optional stopping on $X_t$ or $X_t^2 - t$; It\^o
(Problem~9) produces the two canonical martingales.  Spot the
symmetry, write the martingale, check integrability, apply
optional stopping.

\paragraph{Pattern 3: bound before compute.}
Before computing $\E[T_a]$, ask finite or infinite.  Before
enumerating the LCS table, ask min and max possible answers.  In
the two-children problem, three outcomes with at least one boy
give an instant upper bound of $1/3$.  Bounding first catches
algebra errors and focuses the derivation.

\paragraph{Pattern 4: recurrence to closed form.}
Every DP problem (LCS, knapsack) and random walk problem (coupon
collector, secretary, ruin) has a recurrence.  Seek closed form
when possible (coupon $\to n H_n$; ruin $\to a/N$); accept the
$O(nm)$ table when not (LCS, knapsack).  Knowing which case you
are in, and asking, often decides the interview.

\begin{keyfactbox}[Four meta-patterns]
Every Green-Book-style quant problem yields to one of:
\begin{enumerate}
  \item \textbf{Change of variables} (log, squared, dimensional).
  \item \textbf{Martingale first} (symmetry, optional stopping).
  \item \textbf{Bound before compute} (finite vs infinite, min vs max).
  \item \textbf{Recurrence to closed form} (or an $O(nm)$ table).
\end{enumerate}
If none fires, the problem is a direct consequence of a definition
not yet invoked; re-read the statement for the name of a
distribution, martingale, or recurrence.
\end{keyfactbox}

%% ================================================================
\section{How to practise and what to expect}
\label{sec:prepare}
%% ================================================================

The routine used by the Frankfurt Hedgehogs
\citep{frankfurthedgehogs2025} and other top Prosperity finishers:

\begin{itemize}
  \item Work through the entire Green Book \citeGreenBook{} in
    \emph{writing}, not in your head.  Budget about a month.
  \item For each problem write the five-field template: problem /
    first attempt / what goes wrong / correct / variations.
    Writing the five fields reveals whether you understood or
    merely copied.
  \item Tag each problem with its meta-pattern.  The distribution
    exposes weak areas: 80\% ``change of variables'' and no
    ``martingale'' tells you where to focus.
  \item Do mock interviews out loud.  Whiteboard under time
    pressure is a different skill from written prep; a partner
    asking follow-ups exposes the gaps the real interview will.
\end{itemize}

\begin{prosperitybox}[title={Prosperity skills are interview skills}]
The manual rounds of IMC Prosperity (\cref{ch:manual_rounds}) are
a live, adversarial quant interview: pattern classification of an
unfamiliar product, mental arithmetic under uncertainty,
order-of-magnitude estimation, Kelly sizing.  Top-ten teams
typically treat manual rounds as interview practice, with one
teammate reading and another thinking aloud.  A candidate who has
done five Prosperity manual rounds handles a ``how much would you
pay for this coin-flip game'' question with visible ease.
\end{prosperitybox}

\begin{gsbox}
The Goldman Strats interview pipeline, 2024--2025 cycle
(Engineering, off-cycle):
\begin{itemize}
  \item \textbf{CV screen}, then HireVue behavioural video.
  \item \textbf{First technical phone screen}: 45--60 min, one
    probability and one coding question.  Typical: Monty Hall,
    two-children, derive BS delta; code LCS or knapsack.
  \item \textbf{Superday}: three or four 45-minute technical
    rounds with Strats; one is sometimes a 30-min desk shadow
    followed by a live pricing or hedging question.  At least
    one round tests stochastic calculus at the level of
    \cref{ch:black_scholes}; at least one tests mental arithmetic
    and market intuition (Problem 15 style).
  \item \textbf{Fit / behavioural}: final round with a senior
    Strat or MD.
\end{itemize}
The Green Book \citeGreenBook{} covers $\sim$70\% of the
probability/stochastic content.  The rest splits between numerical
problems (\cref{ch:numerical_methods}), linear algebra, and
desk/product questions (\cref{ch:strats_role}).  The Green Book
alone is necessary but not sufficient.
\end{gsbox}

\begin{historybox}[title={The Green Book}]
Xinfeng Zhou's \emph{A Practical Guide to Quantitative Finance
Interviews} appeared in 2008; Zhou, then a hedge-fund quant,
compiled questions he had asked, been asked, and collected from
colleagues.  The 2020 second edition added ML, more stochastic
calculus, and sharper options coverage.  It remains the de facto
reference for junior quant prep.  A 2026 trading-floor
interviewer almost certainly worked through it, and a candidate
can occasionally turn that to advantage: ``this is close to Green
Book Problem 5.12 but with drift added'' signals preparation
without being obnoxious.
\end{historybox}

%% ================================================================
\section*{Key facts to memorise}
\addcontentsline{toc}{section}{Key facts to memorise}
%% ================================================================

\begin{itemize}
  \item \textbf{Consistency over tricks.}  The five-field template
    applies to every quant interview question; use it in every
    written practice session.
  \item \textbf{Monty Hall $2/3$; two children $1/3$ (ambiguous)
    vs $1/2$ (ordered).}  The difference is what you condition on,
    not the words.
  \item \textbf{100 prisoners win with probability $1 -
    \sum_{51}^{100} 1/\ell \approx 0.312$.}  Cycle-following is
    optimal.
  \item \textbf{Secretary $1/e \approx 0.37$.}  Reject the first
    $n/e$, accept the next best-so-far.
  \item \textbf{Coupon collector $n H_n \approx n \ln n$.}  The
    tail dominates.
  \item \textbf{Brownian $T_a$ is a.s.\ finite, infinite mean,
    $T_a \stackrel{d}{=} a^2 T_1$.}  Dimensional analysis gives the
    scaling.
  \item \textbf{OU half-life $\ln 2 / \theta$.}  Time constant
    $1/\theta$; half-life shorter by $\ln 2$.
  \item \textbf{Gambler's ruin (fair, start $a$, barrier $N$): win
    prob $a/N$, duration $a(N-a)$.}  Two martingales $X_t$ and
    $X_t^2 - t$ give both in two lines.
  \item \textbf{It\^o martingales $W_t^2 - t$ and $e^{W_t - t/2}$.}
    The engines behind Black--Scholes.
  \item \textbf{PCA = max variance = top eigenvector of $\Sigma$;
    rank-1 update $I + uu^\top$ has spectrum $\{1,\dots,1,
    1 + \|u\|^2\}$.}
  \item \textbf{LCS and knapsack are $O(nm)$ DP.}  Knapsack's
    $O(nW)$ is pseudopolynomial, not polynomial.
  \item \textbf{$\Delta = \CDFN(d_1)$.}  Cancellation via
    $S \pdfN(d_1) = Ke^{-rT}\pdfN(d_2)$.  Not $\CDFN(d_2)$, the
    risk-neutral exercise probability.
  \item \textbf{Straddle expected P\&L}: $S_0 \sqrt T \sqrt{2/\pi}
    (\sigma_{\mathrm{real}} - \sigma_{\mathrm{imp}})$.  At 80\% vs
    20\%, $S_0 = 100$, $T = 1$: P\&L $\approx 48$.
  \item \textbf{Meta-patterns.}  Change of variables, martingale
    first, bound before compute, recurrence to closed form: one
    fires on almost every quant interview problem.
\end{itemize}
